\scanpage{17}
\[
 e_{n-1}^{(p)}=-\bigl(\underbrace{A_0(\unsure{\zeta|\xi})}_{a_0}e_0+
 \text{---}+\underbrace{A_{n-1}(\zeta)}_{a_{n-1}}e_{n-1}\bigr)
\]
i.e. on trouve la matrice
\[
 \begin{pmatrix}
 0&0&&0&-a_0\\
 1&0&&0&\unsure{-a_1|a_1}\\
 0&1&\ddots&\vdots&\\
 \bar 0&\bar 0&\cdots&\bar 0&\\
 0&0&\ddots&1&\\
 0&0&&0&-a_{n-1}
 \end{pmatrix}
\]
pour l'application de puiss. \unsure{$p$ième}. Ainsi \deleted{On}
\deleted{\unsure{donnée de}} \inserted{\uline{retrouve que}} la donnée
\deleted{\unsure{du}} d'un sous- \textemdash\textemdash

\diagram{\scriptsize\begin{tikzcd}[ampersand replacement=\&,column sep=small,row sep=large]
 (a_0,-\,a_{n-1})\quad a_0\ \text{inv}
 \arrow[d,<->] \& \subset \&
 (a_0,-,a_{\unsure{n|n-1}})\arrow[d,<->]\\
 \begin{array}{l}\text{ss-gpes }\deleted{\unsure{f}}\text{finis étales de }G_a\end{array}
 \arrow[d,<->] \& \subset \&
 \begin{array}{l}\text{ss-groupes finis loc.}\\\text{libres de rg }p^n\text{ de }G_a\end{array}
 \arrow[d,<->,"?"]\\
 \begin{array}{l}\deleted{E}F\text{ loc. libre }\inserted{\mathrm{rg}\,n}\text{ avec }
 {\deleted{E}F}^{(p)}\simeq\deleted{E}F,\\
 \text{et section }\zeta\text{ de }\deleted{E}F\text{ avec}\\
 \zeta,\zeta^{(p)}\ldots\zeta^{(p^{n-1})}\text{ une base}
 \end{array}
 \arrow[uu,<->,bend right=48] \& \subset \&
 \begin{array}{l}F\text{ loc. libre rg }n\text{ avec}\\
 {\deleted{\unsure{E}}F}^{(p)}\to F\ \text{(pas néc. iso)}\\
 \text{et }\zeta\in\Gamma F\text{ telle que}\\
 \zeta,\zeta^{(p)}\text{---}\zeta^{(p^{n-1})}\text{ base}
 \end{array}
 \arrow[uu,<->,bend right=50,"{\text{\unsure{trivial|trivial\,?}}}"]
\end{tikzcd}}
% archive sheet number: 16
